% ================================================
% EXEMPLE DE SECTION COMPLÈTE AVEC INTÉGRATION GITHUB
% Section 4 : Classes et Objets
% ================================================

\section{Classes et Objets}
\label{sec:classes-objets}

% ================================================
% INTRODUCTION
% ================================================

\subsection{Introduction}

La \textbf{Programmation Orientée Objet} (POO) repose sur le concept fondamental de \emph{classe} et d'\emph{objet}. Une classe est un modèle (ou \textit{blueprint}) qui définit les caractéristiques et comportements d'un type d'objet.

\begin{definition}{Classe}
    Une \textbf{classe} est une structure qui encapsule des données (attributs) et des comportements (méthodes) pour modéliser un concept du monde réel.
\end{definition}

\begin{definition}{Objet}
    Un \textbf{objet} est une instance concrète d'une classe. C'est une entité qui possède un état (valeurs des attributs) et un comportement (méthodes).
\end{definition}

% ================================================
% SYNTAXE
% ================================================

\subsection{Syntaxe d'une Classe}

En Java, une classe se définit avec le mot-clé \java{class} :

\begin{lstlisting}[language=Java, caption=Définition d'une classe Voiture]
public class Voiture {
    // Attributs (variables d'instance)
    private String marque;
    private String modele;
    private int annee;
    private double kilometrage;
    private boolean estDemarree;
    
    // Constructeur par défaut
    public Voiture() {
        this.marque = "Inconnue";
        this.modele = "Inconnu";
        this.annee = 2024;
        this.kilometrage = 0.0;
        this.estDemarree = false;
    }
    
    // Constructeur paramétré
    public Voiture(String marque, String modele, int annee) {
        this.marque = marque;
        this.modele = modele;
        this.annee = annee;
        this.kilometrage = 0.0;
        this.estDemarree = false;
    }
    
    // Méthodes (comportements)
    public void demarrer() {
        if (!estDemarree) {
            estDemarree = true;
            System.out.println(marque + " " + modele + " a démarré");
        }
    }
    
    public void arreter() {
        if (estDemarree) {
            estDemarree = false;
            System.out.println(marque + " " + modele + " s'est arrêtée");
        }
    }
    
    public void rouler(double distance) {
        if (estDemarree) {
            kilometrage += distance;
            System.out.printf("Vous avez parcouru %.1f km%n", distance);
        }
    }
}
\end{lstlisting}

% ================================================
% ENCART GITHUB
% ================================================

\vspace{10pt}
\begin{githubbox}{Ressources GitHub - Section 4}
    \textbf{\faCode\ Exemples de Code}
    \begin{itemize}
        \item \githubexemple{04-Classes-Objets/exemples/Voiture.java}{Voiture.java} - Classe complète avec tests
        \item \githubexemple{04-Classes-Objets/exemples/CompteBancaire.java}{CompteBancaire.java} - Gestion de compte bancaire
        \item \href{\repourl/blob/main/04-Classes-Objets/exemples/screenshots/output-examples.txt}{\faTerminal\ Sorties d'exécution}
    \end{itemize}
    
    \vspace{5pt}
    
    \textbf{\faPencilAlt\ Exercices Pratiques}
    \begin{itemize}
        \item \githubexercice{04-Classes-Objets/exercices/exercice1.md}{Énoncé : Classe Étudiant}
        \item \githubexemple{04-Classes-Objets/solutions/Exercice1Solution.java}{Solution complète}
    \end{itemize}
    
    \vspace{5pt}
    
    \textit{Accès direct :} \href{\repourl/tree/main/04-Classes-Objets}{\texttt{04-Classes-Objets/}}
\end{githubbox}

% ================================================
% UTILISATION
% ================================================

\subsection{Création et Utilisation d'Objets}

Pour utiliser une classe, on crée des \textbf{instances} (objets) :

\begin{lstlisting}[language=Java, caption=Utilisation de la classe Voiture]
public class TestVoiture {
    public static void main(String[] args) {
        // Création d'objets (instanciation)
        Voiture voiture1 = new Voiture();
        Voiture voiture2 = new Voiture("Peugeot", "308", 2023);
        
        // Utilisation des méthodes
        voiture2.demarrer();
        voiture2.rouler(50.5);
        voiture2.arreter();
    }
}
\end{lstlisting}

\begin{remarque}
    L'opérateur \java{new} alloue de la mémoire pour l'objet et appelle le constructeur correspondant.
\end{remarque}

% ================================================
% ENCAPSULATION
% ================================================

\subsection{Encapsulation}

L'\textbf{encapsulation} est un principe fondamental de la POO qui consiste \`a :
\begin{itemize}
    \item Rendre les attributs \java{private}
    \item Fournir des méthodes \java{public} pour y accéder (\textit{getters} et \textit{setters})
\end{itemize}

\begin{exemple}{Getters et Setters}
\begin{lstlisting}[language=Java]
public class Voiture {
    private String marque;
    
    // Getter
    public String getMarque() {
        return marque;
    }
    
    // Setter avec validation
    public void setMarque(String marque) {
        if (marque != null && !marque.isEmpty()) {
            this.marque = marque;
        } else {
            throw new IllegalArgumentException("Marque invalide");
        }
    }
}
\end{lstlisting}
\end{exemple}

% ================================================
% EXERCICES
% ================================================

\subsection{Exercices}

\begin{exercice}{1 - Classe Étudiant}
    Créer une classe \javaclass{Etudiant} avec les attributs suivants :
    \begin{itemize}
        \item \java{matricule} (String)
        \item \java{nom} et \java{prenom} (String)
        \item \java{moyenne} (double)
        \item \java{nbAbsences} (int)
    \end{itemize}
    
    Implémenter :
    \begin{itemize}
        \item Constructeurs (par défaut et paramétré)
        \item Getters et setters avec validation
        \item Méthode \javamethod{estAdmis} (retourne \java{true} si moyenne $\geq$ 10)
        \item Méthode \javamethod{afficher} (affiche toutes les informations)
    \end{itemize}
    
    \vspace{5pt}
    
    \textbf{\faGithub\ Solution :} \githubexemple{04-Classes-Objets/solutions/Exercice1Solution.java}{Voir sur GitHub}
\end{exercice}

% ================================================
% RÉSUMÉ
% ================================================

\subsection{Points Clés}

\begin{itemize}
    \item Une \textbf{classe} est un modèle, un \textbf{objet} est une instance
    \item Les \textbf{attributs} représentent l'état, les \textbf{méthodes} le comportement
    \item L'\textbf{encapsulation} protège les données
    \item Les \textbf{constructeurs} initialisent les objets
    \item Le mot-clé \java{this} référence l'objet courant
\end{itemize}

% ================================================
% FIN DE LA SECTION
% ================================================
